\documentclass[aspectratio=169,10pt]{beamer}

% ==============================================================================
% CẤU HÌNH GÓI LỆNH VÀ FONT CHỮ (XeLaTeX)
% ==============================================================================
\usepackage{fontspec}
\setsansfont{Segoe UI}[
  BoldFont = Segoe UI Bold,
  ItalicFont = Segoe UI Italic,
  BoldItalicFont = Segoe UI Bold Italic
]
\setmonofont{Consolas}

\usepackage{amsmath,amssymb,amsfonts}
\usepackage{bm}
\usepackage{booktabs}
\usepackage{tabularx}
\usepackage{multirow}
\usepackage{multicol}
\usepackage{tikz}
\usetikzlibrary{calc,shapes.geometric,arrows,arrows.meta,positioning,fit}
\usepackage{graphicx}
\graphicspath{{fig/}{paper/fig/}{paper/}{./}}

% ==============================================================================
% BẢNG MÀU THIẾT KẾ HIỆN ĐẠI (CHUẨN NHẬN DIỆN PTIT)
% ==============================================================================
\definecolor{ptitblue}{RGB}{15, 44, 89}      % Xanh Navy chủ đạo (#0F2C59)
\definecolor{ptitred}{RGB}{186, 37, 37}      % Đỏ PTIT điểm nhấn (#BA2525)
\definecolor{ptitdark}{RGB}{30, 41, 59}      % Xám than văn bản (#1E293B)
\definecolor{ptitlight}{RGB}{241, 245, 249}  % Nền xám nhạt (#F1F5F9)
\definecolor{ptitaccent}{RGB}{2, 132, 199}   % Xanh biển phụ trợ (#0284C7)
\definecolor{ptitgreen}{RGB}{22, 163, 74}    % Xanh lá kết quả (#16A34A)

% Tùy biến giao diện Beamer phẳng, tinh tế
\setbeamercolor{normal text}{fg=ptitdark,bg=white}
\setbeamercolor{structure}{fg=ptitblue}
\setbeamercolor{frametitle}{fg=white,bg=ptitblue}
\setbeamercolor{title}{fg=white,bg=ptitblue}
\setbeamercolor{subtitle}{fg=white!85,bg=ptitblue}
\setbeamercolor{block title}{fg=white,bg=ptitblue}
\setbeamercolor{block body}{fg=ptitdark,bg=ptitlight}
\setbeamercolor{block title alerted}{fg=white,bg=ptitred}
\setbeamercolor{block body alerted}{fg=ptitdark,bg=ptitred!8}
\setbeamercolor{block title example}{fg=white,bg=ptitaccent}
\setbeamercolor{block body example}{fg=ptitdark,bg=ptitaccent!8}

\setbeamertemplate{navigation symbols}{}
\setbeamertemplate{items}[circle]
\setbeamerfont{frametitle}{series=\bfseries,size=\large}
\setbeamerfont{title}{series=\bfseries,size=\Large}

% Footline hiển thị trang và thông tin đề tài
\setbeamertemplate{footline}{
  \leavevmode%
  \hbox{%
  \begin{beamercolorbox}[wd=.35\paperwidth,ht=2.4ex,dp=0.8ex,left,leftskip=0.8cm]{normal text}%
    \usebeamerfont{author in head/foot}\scriptsize\textbf{Lê Việt Anh} -- PTIT (24-SV-CNTT-2026)
  \end{beamercolorbox}%
  \begin{beamercolorbox}[wd=.45\paperwidth,ht=2.4ex,dp=0.8ex,center]{normal text}%
    \usebeamerfont{title in head/foot}\scriptsize\textbf{Báo cáo Đề tài Nghiên cứu Khoa học Sinh viên}
  \end{beamercolorbox}%
  \begin{beamercolorbox}[wd=.20\paperwidth,ht=2.4ex,dp=0.8ex,right,rightskip=0.8cm]{normal text}%
    \scriptsize\textbf{\insertframenumber{} / \inserttotalframenumber}
  \end{beamercolorbox}}%
  \vskip0pt%
}

% Định dạng tiêu đề slide
\setbeamertemplate{frametitle}{
  \begin{beamercolorbox}[wd=\paperwidth,ht=3.0ex,dp=1.0ex,leftskip=0.6cm]{frametitle}
    \insertframetitle
  \end{beamercolorbox}
}

% Macro tạo nhãn badge nổi bật
\newcommand{\badge}[2]{\colorbox{#1!15}{\textcolor{#1}{\textbf{\scriptsize #2}}}}
\newcommand{\metricbox}[3]{
  \begin{beamercolorbox}[rounded=true,shadow=false,wd=\linewidth,sep=3pt]{block body}
    \centering
    {\footnotesize \textbf{#1}}\\[1pt]
    {\large \textbf{\textcolor{ptitred}{#2}}}\\[1pt]
    {\tiny \textcolor{black!60}{#3}}
  \end{beamercolorbox}
}

% ==============================================================================
% THÔNG TIN TIÊU ĐỀ BÁO CÁO
% ==============================================================================
\title{NGHIÊN CỨU, ĐỀ XUẤT MÔ HÌNH HỌC SÂU\\TRONG DỰ ĐOÁN TRẠNG THÁI TRONG MẠNG LƯỚI\\GIAO THÔNG ĐƯỜNG BỘ}
\subtitle{Báo cáo Nghiệm thu Đề tài Nghiên cứu Khoa học Cấp Học viện -- Năm 2026}
\author{\textbf{Sinh viên thực hiện:} Lê Việt Anh (B23DCKH002)\\[2pt]
\textbf{Giảng viên hướng dẫn:} PGS. TS. Nguyễn Trọng Khánh}
\institute{Khoa Công nghệ Thông tin 1 -- Học viện Công nghệ Bưu chính Viễn thông}
\date{Hà Nội, Tháng 10 năm 2026}

\begin{document}

% ------------------------------------------------------------------------------
% SLIDE 1: TRANG TIÊU ĐỀ
% ------------------------------------------------------------------------------
\begin{frame}[plain]
  \begin{tikzpicture}[remember picture,overlay]
    \fill[ptitblue] (current page.south west) rectangle (current page.north east);
    \fill[ptitred] ($(current page.north west)+(0,-1.6cm)$) rectangle ($(current page.north east)+(0,-1.7cm)$);
    \node[anchor=north west] at ($(current page.north west)+(0.8cm,-0.3cm)$) {
      \includegraphics[height=1.0cm]{image2.png}
    };
    \node[anchor=north west, text=white, align=left] at ($(current page.north west)+(2.2cm,-0.35cm)$) {
      \small\textbf{HỌC VIỆN CÔNG NGHỆ BƯU CHÍNH VIỄN THÔNG}\\
      \scriptsize\textbf{KHOA CÔNG NGHỆ THÔNG TIN 1}
    };
    \node[anchor=north east, text=white] at ($(current page.north east)+(-0.8cm,-0.35cm)$) {
      \scriptsize\badge{white}{\textcolor{ptitred}{MÃ ĐỀ TÀI: 24-SV-CNTT-2026}}
    };
    \node[text width=0.88\paperwidth, align=center, text=white] at ($(current page.center)+(0, 0.6cm)$) {
      {\Large \textbf{NGHIÊN CỨU, ĐỀ XUẤT MÔ HÌNH HỌC SÂU}}\\[4pt]
      {\Large \textbf{TRONG DỰ ĐOÁN TRẠNG THÁI TRONG MẠNG LƯỚI}}\\[4pt]
      {\Large \textbf{GIAO THÔNG ĐƯỜNG BỘ}}\\[8pt]
      {\small\textcolor{white!85}{Khung mô hình hai giai đoạn từ camera giám sát và kiểm soát lan truyền sai số}}
    };
    \node[text width=0.88\paperwidth, text=white] at ($(current page.south)+(0, 1.4cm)$) {
      \begin{columns}[T]
        \begin{column}{0.48\textwidth}
          \small \textbf{Giảng viên hướng dẫn:}\\
          \textbf{PGS. TS. Nguyễn Trọng Khánh}
        \end{column}
        \begin{column}{0.48\textwidth}
          \raggedleft
          \small \textbf{Sinh viên thực hiện:}\\
          \textbf{Lê Việt Anh} -- B23DCKH002
        \end{column}
      \end{columns}
      \vskip 4pt
      \centering
      {\scriptsize \textcolor{white!70}{Hà Nội -- Năm 2026}}
    };
  \end{tikzpicture}
\end{frame}

% ------------------------------------------------------------------------------
% SLIDE 2: BỐI CẢNH & ĐẶT VẤN ĐỀ
% ------------------------------------------------------------------------------
\begin{frame}{1. Bối cảnh và Đặt vấn đề nghiên cứu}
  \begin{columns}[T]
    \begin{column}{0.50\textwidth}
      \begin{block}{\textbf{Thực trạng giao thông đô thị Việt Nam}}
        \footnotesize
        \begin{itemize}
          \item \textbf{Dòng xe hỗn hợp đặc thù:} Xe máy chiếm $>80\%$ lưu lượng, di chuyển linh hoạt, không tuân thủ làn cứng.
          \item \textbf{Hạn chế của cảm biến vòng từ (Inductive Loops):}
            \begin{itemize}
              \scriptsize
              \item Không phân biệt được nhiều xe máy di chuyển song song.
              \item Chi phí xẻ đường, lắp đặt và bảo trì định kỳ rất lớn.
            \end{itemize}
          \item \textbf{Cơ hội từ hạ tầng camera giám sát:} Hệ thống camera giao thông sẵn có tại TP.HCM cung cấp dữ liệu hình ảnh liên tục với chi phí thấp.
        \end{itemize}
      \end{block}
      \vskip 2pt
      \badge{ptitblue}{Mục tiêu:} \footnotesize Khai thác triệt để ảnh camera giám sát để ước lượng và dự báo lưu lượng dòng xe hỗn hợp toàn đô thị.
    \end{column}

    \begin{column}{0.47\textwidth}
      \centering
      \includegraphics[width=\linewidth,height=4.3cm,keepaspectratio]{samples_pictures_get_from_api.png}
      \vskip 2pt
      {\tiny \textbf{Hình 1: Dữ liệu ảnh thực tế thu nhận từ hệ thống camera giao thông TP.HCM}}
      \vskip 2pt
      \begin{beamercolorbox}[rounded=true,wd=\linewidth,sep=3pt]{block body alerted}
        \scriptsize
        \textbf{Thách thức thị giác:} Mật độ phương tiện dày đặc, che khuất lẫn nhau nghiêm trọng, thay đổi góc quay và điều kiện ánh sáng ngày/đêm.
      \end{beamercolorbox}
    \end{column}
  \end{columns}
\end{frame}

% ------------------------------------------------------------------------------
% SLIDE 3: HIỆN TƯỢNG LAN TRUYỀN SAI SỐ
% ------------------------------------------------------------------------------
\begin{frame}{2. Thách thức khoa học: Hiện tượng Lan truyền sai số (Error Propagation)}
  \begin{columns}[T]
    \begin{column}{0.48\textwidth}
      \begin{alertblock}{\textbf{Cơ chế phát sinh \& khuếch tán sai số}}
        \footnotesize
        \begin{itemize}
          \item \textbf{Sai số ước lượng thị giác (Giai đoạn 1):} Do góc nghiêng camera và che khuất, số đếm luôn tồn tại sai số $\boldsymbol{\eta}_i(t)$:
            \begin{equation*}
              \mathbf{x}_i(t) = \mathbf{c}_i(t) + \boldsymbol{\eta}_i(t)
            \end{equation*}
          \item \textbf{Lan truyền theo không gian:} Phép tích chập đồ thị phổ ChebNet khuếch tán sai số $\boldsymbol{\eta}$ sang các nút lân cận trên mạng lưới đường bộ.
          \item \textbf{Tích lũy theo thời gian:} Khi dự báo đa bước ($H=6$, 30 phút), sai số ngắn hạn làm chệch hướng dự báo dài hạn.
        \end{itemize}
      \end{alertblock}
    \end{column}

    \begin{column}{0.48\textwidth}
      \begin{block}{\textbf{Chuỗi lan truyền sai số nối tầng trong Two-Stage Pipeline}}
        \centering
        \vskip 4pt
        \begin{tikzpicture}[
          >=latex, font=\scriptsize,
          box/.style={draw=ptitblue, thick, rounded corners=2pt, fill=white, text centered, inner sep=3pt},
          err/.style={draw=ptitred, thick, rounded corners=2pt, fill=ptitred!10, text centered, inner sep=3pt},
          arr/.style={draw=ptitblue, -{Stealth[scale=1.0]}, thick}
        ]
          \node[box, text width=3.3cm] (raw) {Ảnh camera tuyến đường $\mathbf{I}_t$};
          \node[box, text width=3.3cm, below=0.35cm of raw] (stg1) {Khâu thị giác $f_{\theta_1}$ (Đếm trực tiếp)};
          \node[err, text width=3.5cm, below=0.35cm of stg1] (noise) {Lưu lượng chứa sai số\\ $\mathbf{x}_i(t) = \mathbf{c}_i(t) + \boldsymbol{\eta}_i(t)$};
          \node[box, text width=3.3cm, below=0.35cm of noise] (stg2) {Khâu đồ thị $g_{\theta_2}$ (Dự báo mạng lưới)};
          \node[err, text width=3.5cm, below=0.35cm of stg2] (out) {Sai số lan truyền \& tích lũy\\ $\hat{\mathbf{Y}}_{t+1:t+H}$ bị suy thoái};

          \draw[arr] (raw) -- (stg1);
          \draw[arr] (stg1) -- (noise);
          \draw[arr] (noise) -- (stg2);
          \draw[arr] (stg2) -- (out);
        \end{tikzpicture}
      \end{block}
    \end{column}
  \end{columns}
\end{frame}

% ------------------------------------------------------------------------------
% SLIDE 4: MỤC TIÊU & ĐÓNG GÓP CHÍNH
% ------------------------------------------------------------------------------
\begin{frame}{3. Mục tiêu nghiên cứu và Đóng góp chính của đề tài}
  \begin{columns}[T]
    \begin{column}{0.48\textwidth}
      \textbf{\textcolor{ptitblue}{Bốn mục tiêu cụ thể:}}
      \begin{enumerate}
        \footnotesize
        \item \textbf{Bộ dữ liệu IC4SD-Traffic-HCM:} 608 nút camera giám sát TP.HCM, chuỗi thời gian 86 ngày liên tục, chuẩn hóa 2 lớp phương tiện.
        \item \textbf{Hồi quy đếm xe trực tiếp (Giai đoạn 1):} Loại bỏ hoàn toàn hộp bao và thuật toán NMS, giải quyết triệt để bài toán che khuất xe máy.
        \item \textbf{Mô hình đồ thị TA-STGCN (Giai đoạn 2):} Tích hợp cơ chế chú ý thời gian đa đầu (MHTA) nhằm chặn đứng lan truyền sai số.
        \item \textbf{Thiết kế hệ thống phân tán nối tầng:} Tối ưu hóa tính độc lập giữa cảm nhận thị giác biên và dự báo trung tâm.
      \end{enumerate}
    \end{column}

    \begin{column}{0.48\textwidth}
      \textbf{\textcolor{ptitred}{Ba đóng góp khoa học then chốt:}}
      \vskip 4pt
      \metricbox{Bộ dữ liệu thực tế quy mô đô thị}{608 Camera / 86 Ngày}{Bảo toàn 2 lớp ô tô \& xe máy, lấy mẫu 5 phút chuẩn tắc}
      \vskip 3pt
      \metricbox{Hiệu năng kìm hãm sai số vượt trội}{Cohen's $d = +5.84$}{Vượt trội GraphWaveNet, nhẹ hơn 70.4\% tham số so với STAEformer}
      \vskip 3pt
      \metricbox{Tốc độ suy luận đồ thị toàn mạng}{1.23 ms / Mạng lưới}{Bộ nhớ huấn luyện đỉnh 8.6 GB VRAM, đáp ứng thời gian thực}
    \end{column}
  \end{columns}
\end{frame}

% ------------------------------------------------------------------------------
% SLIDE 5: BỘ DỮ LIỆU CHUẨN IC4SD-TRAFFIC-HCM
% ------------------------------------------------------------------------------
\begin{frame}{4. Bộ dữ liệu chuẩn mực IC4SD-Traffic-HCM}
  \begin{columns}[T]
    \begin{column}{0.52\textwidth}
      \textbf{\textcolor{ptitblue}{Cấu trúc 3 cấu phần chuẩn hóa:}}
      \begin{itemize}
        \footnotesize
        \item \textbf{IC4SD-Traffic-HCM-Vision (Tập ảnh Giai đoạn 1):}
          \begin{itemize}
            \scriptsize
            \item 5.012 ảnh tĩnh độ phân giải cao từ camera giao thông.
            \item 3 thẩm định viên gán nhãn: Tương quan $R^2 \ge 0.968$.
            \item Chia không gian (Spatial Disjoint): 4.010 Train (80\%) / 501 Val (10\%) / 501 Test (10\%).
          \end{itemize}
        \item \textbf{IC4SD-Traffic-HCM-Graph (Chuỗi Giai đoạn 2):}
          \begin{itemize}
            \scriptsize
            \item Chuỗi 86 ngày liên tục ($N=608$ nút trạm).
            \item 24.768 bước quan sát ($\Delta t = 5$ ph). Lọc trượt kép nhân quả.
            \item Chia theo thời gian: 80\% Train / 10\% Val / 10\% Test.
          \end{itemize}
        \item \textbf{Tô-pô đồ thị OpenStreetMap (OSM):}
          \begin{itemize}
            \scriptsize
            \item $N = 608$ nút trạm, ma trận khoảng cách $W_{ij} = \exp(-d_{ij}^2/\sigma^2)$.
          \end{itemize}
      \end{itemize}
    \end{column}

    \begin{column}{0.46\textwidth}
      \begin{block}{\textbf{Thông số kỹ thuật bộ dữ liệu đồ thị}}
        \centering
        \fontsize{6.5pt}{7.5pt}\selectfont
        \renewcommand{\arraystretch}{1.05}
        \setlength{\tabcolsep}{3pt}
        \begin{tabular}{lc}
          \toprule
          \textbf{Thông số kỹ thuật} & \textbf{Giá trị} \\
          \midrule
          Số trạm camera giám sát ($N$) & \textbf{608} nút \\
          Phân loại phương tiện theo dõi & \textbf{2} (Ô tô, Xe máy) \\
          Độ dài chuỗi thời gian & \textbf{86} ngày \\
          Cửa sổ quan sát lịch sử ($T_{\text{in}}$) & \textbf{24} bước (120 phút) \\
          Tầm nhìn dự báo ($H$) & \textbf{6} bước (30 phút) \\
          Tổng số bước thời gian & \textbf{24.768} bước \\
          Mẫu trượt (Train / Val / Test) & \textbf{19.785 / 2.472 / 2.472} \\
          Tỷ lệ phân chia Chronological & \textbf{80\% / 10\% / 10\%} \\
          \bottomrule
        \end{tabular}
      \end{block}
      \vskip 3pt
      \centering
      {\scriptsize \textbf{\textcolor{ptitgreen}{Bộ dữ liệu quy mô lớn đầu tiên cho dòng xe hỗn hợp TP.HCM}}}
    \end{column}
  \end{columns}
\end{frame}

% ------------------------------------------------------------------------------
% SLIDE 6: TỔNG QUAN KHUNG MÔ HÌNH HAI GIAI ĐOẠN
% ------------------------------------------------------------------------------
\begin{frame}{5. Tổng quan Khung mô hình hai giai đoạn (Two-Stage Pipeline)}
  \centering
  \begin{tikzpicture}[
    >=latex, font=\scriptsize,
    proc/.style={draw=ptitblue, thick, fill=white, rounded corners=3pt, text centered, inner sep=4pt},
    data/.style={draw=ptitdark, fill=ptitlight, rounded corners=3pt, text centered, inner sep=4pt},
    res/.style={draw=ptitred, thick, fill=ptitred!8, rounded corners=3pt, text centered, inner sep=4pt},
    arr/.style={draw=ptitblue, -{Stealth[scale=1.1]}, thick}
  ]
    % Nodes
    \node[data, text width=2.5cm] (cam) at (0, 1.8) {
      \textbf{Camera Giám sát}\\
      $\mathbf{I}_i(t) \in \mathbb{R}^{3 \times 224 \times 224}$\\
      \textit{Ảnh tĩnh camera tuyến đường}
    };

    \node[proc, text width=3.3cm] (stg1) at (4.0, 1.8) {
      \textbf{GIAI ĐOẠN 1 (Vision)}\\
      EfficientNet-B5 + Bottleneck MLP\\
      \textit{Hồi quy số đếm trực tiếp}
    };

    \node[data, text width=2.8cm] (count) at (8.5, 1.8) {
      \textbf{Vector số đếm nút}\\
      $[\hat{c}_{\text{car}}, \hat{c}_{\text{bike}}]^T \in \mathbb{R}^2_{\ge 0}$\\
      \textcolor{ptitgreen}{\textbf{Lưu lượng tức thời}}
    };

    \node[data, text width=2.8cm] (mat) at (8.5, -0.8) {
      \textbf{Tensor mạng lưới}\\
      $\mathbf{X} \in \mathbb{R}^{B \times 24 \times 608 \times 5}$\\
      \textit{2 kênh xe + 3 kênh thời gian}
    };

    \node[proc, text width=3.3cm] (stg2) at (4.0, -0.8) {
      \textbf{GIAI ĐOẠN 2 (Graph)}\\
      \textbf{Mô hình TA-STGCN}\\
      ChebNet + MHTA + GLU ($C=80$)
    };

    \node[res, text width=2.5cm] (out) at (0, -0.8) {
      \textbf{Dự báo đa bước}\\
      $\hat{\mathbf{Y}} \in \mathbb{R}^{B \times 6 \times 608 \times 2}$\\
      \textit{Tầm nhìn nhìn xa 30 phút}
    };

    % Arrows
    \draw[arr] (cam) -- (stg1);
    \draw[arr] (stg1) -- (count);
    \draw[arr] (count) -- node[right, font=\tiny] {Ghép 608 nút + Time} (mat);
    \draw[arr] (mat) -- (stg2);
    \draw[arr] (stg2) -- (out);

    % Boundary note
    \draw[ptitred, dashed, thick] (-1.5, 0.5) -- (10.5, 0.5);
    \node[anchor=west, text=ptitred, font=\tiny\bfseries] at (-1.4, 0.7) {Nguyên tắc phân tách độc lập (Decoupling): Phân hệ thị giác và đồ thị liên kết lỏng qua vector số đếm};
  \end{tikzpicture}
\end{frame}

% ------------------------------------------------------------------------------
% SLIDE 7: GIAI ĐOẠN 1: HỒI QUY ĐẾM XE TRỰC TIẾP
% ------------------------------------------------------------------------------
\begin{frame}{6. Giai đoạn 1: Hồi quy đếm xe trực tiếp từ ảnh camera}
  \begin{columns}[T]
    \begin{column}{0.51\textwidth}
      \begin{block}{\textbf{Nguyên lý hồi quy đếm xe trực tiếp}}
        \scriptsize
        \begin{itemize}
          \item \textbf{Loại bỏ Bounding Box \& NMS:} Tránh bỏ sót xe máy che khuất nhau, triệt tiêu độ phức tạp $\mathcal{O}(M^2)$.
          \item \textbf{EfficientNet-B5 Backbone:} Trích xuất đặc trưng sâu bằng MBConv kết hợp SE attention $\to \mathbf{f} \in \mathbb{R}^{2048}$.
          \item \textbf{Bottleneck MLP + Huber Loss:} Nén $2048 \to 128 \to 2$, tối ưu hóa bền vững trước ảnh lóa/ngập nước.
        \end{itemize}
      \end{block}
      \vskip 1pt
      \begin{block}{\textbf{Đối sánh hiệu năng các Backbone thị giác}}
        \centering
        \resizebox{\linewidth}{!}{%
        \begin{tabular}{lccccc}
          \toprule
          \textbf{Mô hình} & \textbf{Car MAE} $\downarrow$ & \textbf{Bike MAE} $\downarrow$ & \textbf{Macro MAE} $\downarrow$ & \textbf{GFLOPs} $\downarrow$ & \textbf{FPS} $\uparrow$ \\
          \midrule
          YOLOv8x-Crowd & 3.42 & 6.18 & 4.800 & 258.0 & 18.5 \\
          MobileNetV3-L & 2.81 & 5.09 & 3.950 & \textbf{0.44} & \textbf{142.8} \\
          ResNet-50 & 2.45 & 4.62 & 3.535 & 8.2 & 84.2 \\
          ConvNeXt-Base & 2.22 & 4.28 & 3.250 & 30.8 & 42.6 \\
          Swin-Base & 2.20 & 4.25 & 3.225 & 31.0 & 38.2 \\
          \midrule
          \textbf{EffNet-B5 (Đề xuất)} & \textbf{2.14} & \textbf{4.18} & \textbf{3.160} & \textbf{20.6} & \textbf{52.4} \\
          \bottomrule
        \end{tabular}%
        }
      \end{block}
    \end{column}

    \begin{column}{0.47\textwidth}
      \centering
      \includegraphics[width=\linewidth,height=3.8cm,keepaspectratio]{vision_explainability_gradcam.png}
      \vskip 2pt
      {\tiny \textbf{Hình 2: Giải thích tính chú ý của mạng thị giác bằng Grad-CAM++}}
      \vskip 2pt
      \begin{beamercolorbox}[rounded=true,wd=\linewidth,sep=3pt]{block body example}
        \scriptsize
        \textbf{Phân tích trực quan Grad-CAM++:} EfficientNet-B5 tập trung vùng kích hoạt chuẩn xác vào thân xe ô tô và cụm đầu xe máy, không bị phân tán bởi cảnh nền đường xá.
      \end{beamercolorbox}
    \end{column}
  \end{columns}
\end{frame}

% ------------------------------------------------------------------------------
% SLIDE 8: GIAI ĐOẠN 2: KIẾN TRÚC MÔ HÌNH ĐỀ XUẤT TA-STGCN
% ------------------------------------------------------------------------------
\begin{frame}{7. Giai đoạn 2: Kiến trúc mô hình đề xuất TA-STGCN}
  \begin{columns}[T]
    \begin{column}{0.48\textwidth}
      \textbf{\textcolor{ptitblue}{Cấu trúc khối xử lý không -- thời gian:}}
      \begin{enumerate}
        \footnotesize
        \item \textbf{Tích chập không gian ChebNet ($K=3$):}
          \begin{equation*}
            g_\theta \star_{\mathcal{G}} \mathbf{x} = \sum_{k=0}^{K-1} \theta_k T_k(\tilde{\mathbf{L}}) \mathbf{x}
          \end{equation*}
          Khai thác đồ thị vật lý OpenStreetMap với chi phí tuyến tính $\mathcal{O}(K|\mathcal{E}|)$.
        \item \textbf{Tích chập thời gian có cổng (1D GLU):}
          Kernel $K_t=3, \text{pad}=1$ bảo toàn $T_{\text{in}}=24$, xử lý song song toàn bộ chuỗi thời gian.
        \item \textbf{Cơ chế Chú ý Thời gian Đa đầu (MHTA):}
          Bố trí dạng sandwich ($h=4, d_k=20$) tái cân bằng trọng số giữa các mốc thời gian lịch sử.
      \end{enumerate}
      \vskip 2pt
      \badge{ptitblue}{Tham số: \textbf{453.052}} \quad
      \badge{ptitgreen}{VRAM Train: \textbf{8.6 GB}} \quad
      \badge{ptitaccent}{Suy luận: \textbf{1.23 ms}}
    \end{column}

    \begin{column}{0.50\textwidth}
      \centering
      \includegraphics[width=\linewidth,height=4.6cm,keepaspectratio]{ta_stgcn_architecture.pdf}
      \vskip 2pt
      {\tiny \textbf{Hình 3: Sơ đồ khối kiến trúc mạng đề xuất TA-STGCN}}
      \vskip 3pt
      \begin{beamercolorbox}[rounded=true,wd=\linewidth,sep=3pt]{block body}
        \scriptsize
        \textbf{Cấu trúc Sandwich:} ST-Conv 1 ($C=80$) $\to$ Multi-Head Temporal Attention ($h=4$) $\to$ ST-Conv 2 ($C=80$) $\to$ Tầng chiếu 1D Conv ngõ ra.
      \end{beamercolorbox}
    \end{column}
  \end{columns}
\end{frame}

% ------------------------------------------------------------------------------
% SLIDE 9: CƠ CHẾ KHỬ NHIỄU CỦA TEMPORAL ATTENTION
% ------------------------------------------------------------------------------
\begin{frame}{8. Cơ chế khử nhiễu của Temporal Self-Attention}
  \begin{columns}[T]
    \begin{column}{0.48\textwidth}
      \begin{alertblock}{\textbf{Nguyên lý toán học: Bộ lọc thông thấp thích ứng}}
        \scriptsize
        Khi đầu vào chứa sai số $\boldsymbol{\eta} \sim (0, \sigma_\eta^2)$, Attention tính tổ hợp lồi:\\
        \centering
        $\operatorname{Var}\left(\sum_{t=1}^{T_{\text{in}}} \alpha_t \eta_t\right) = \sigma_\eta^2 \sum_{t=1}^{T_{\text{in}}} \alpha_t^2 = \frac{\sigma_\eta^2}{T_{\text{eff}}}$ với $T_{\text{eff}} = \frac{1}{\|\boldsymbol{\alpha}\|_2^2}$\\[2pt]
        \raggedright
        $\implies$ \textbf{Phương sai sai số giảm tỷ lệ nghịch với $T_{\text{eff}}$}!
      \end{alertblock}

      \begin{block}{\textbf{Định lượng chú ý qua các khung giờ (Bảng 5.11)}}
        \scriptsize
        \begin{itemize}
          \item \textbf{Đêm khuya (00:00--05:00):} $T_{\text{eff}} = 16.2 \pm 0.8$, Entropy $2.84 \implies$ Phân bổ đều để triệt tiêu nhiễu ánh sáng tối.
          \item \textbf{Cao điểm chiều (17:00--19:00):} $T_{\text{eff}} = 8.4 \pm 0.5$, Entropy $2.15 \implies$ Tập trung bước cận kề nắm bắt ùn tắc.
        \end{itemize}
      \end{block}
    \end{column}

    \begin{column}{0.50\textwidth}
      \centering
      \includegraphics[width=\linewidth,height=3.6cm,keepaspectratio]{attention_comparison_evening_vs_noon.pdf}
      \vskip 2pt
      {\tiny \textbf{Hình 4: Ma trận trọng số chú ý thời gian: Giờ cao điểm chiều vs Trưa bình thường}}
      \vskip 2pt
      \begin{beamercolorbox}[rounded=true,wd=\linewidth,sep=3pt]{block body example}
        \scriptsize
        \textbf{Quan sát thực nghiệm:} Mạng tự động điều chỉnh cơ chế chú ý: tập trung cao độ vào các bước quan sát tức thời khi tắc đường và trải rộng làm mượt khi lưu lượng ổn định.
      \end{beamercolorbox}
    \end{column}
  \end{columns}
\end{frame}

% ------------------------------------------------------------------------------
% SLIDE 10: KẾT QUẢ THỰC NGHIỆM ĐỒ THỊ (SOTA BENCHMARK)
% ------------------------------------------------------------------------------
\begin{frame}{9. Kết quả thực nghiệm dự báo trên Benchmark đồ thị}
  \begin{block}{\textbf{Bảng đối sánh hiệu năng dự báo đa bước ($H=6$, 30 phút) trên toàn mạng 608 nút (Trung bình 5 random seeds)}}
    \centering
    \fontsize{6.2pt}{7.2pt}\selectfont
    \renewcommand{\arraystretch}{1.03}
    \setlength{\tabcolsep}{3.0pt}
    \begin{tabular}{lcccccccc}
      \toprule
      \textbf{Mô hình} & \textbf{MAE $t+1$} $\downarrow$ & \textbf{MAE $t+3$} $\downarrow$ & \textbf{MAE $t+6$} $\downarrow$ & \textbf{Overall MAE} $\downarrow$ & \textbf{RMSE} $\downarrow$ & \textbf{MAPE} (\%) $\downarrow$ & \textbf{Params} & \textbf{VRAM (MB)} \\
      \midrule
      STGCN (2018) & 2.9461 & 3.3248 & 3.6521 & 3.3077 & 5.3102 & 14.38 & 375K & 4.838 \\
      GraphWaveNet (2019) & 2.8715 & 3.2234 & 3.5800 & 3.2250 & 5.2910 & 14.25 & 312K & 5.120 \\
      ASTGCN (2019) & 2.9840 & 3.3615 & 3.7102 & 3.3519 & 5.4210 & 14.80 & 410K & 5.640 \\
      DSTAGNN (2022) & 2.8950 & 3.2510 & 3.5940 & 3.2467 & 5.2815 & 14.20 & 890K & 7.210 \\
      MegaCRN (2023) & 2.8820 & 3.2390 & 3.5780 & 3.2330 & 5.2750 & 14.18 & 1.25M & 8.940 \\
      STAEformer (2023) & \textbf{2.8410} & 3.1950 & 3.5120 & 3.1827 & \textbf{5.2490} & 14.08 & 1.53M & 11.520 \\
      iTransformer (2024) & 2.9120 & 3.2780 & 3.6210 & 3.2703 & 5.3010 & 14.31 & \textbf{265K} & \textbf{265} \\
      \midrule
      \textbf{TA-STGCN (Đề xuất)} & 2.8524 & \textbf{3.1892} & \textbf{3.4947} & \textbf{3.1788} & 5.2536 & \textbf{13.82} & \textbf{453K} & \textbf{8.600} \\
      \bottomrule
    \end{tabular}
  \end{block}

  \begin{columns}[T]
    \begin{column}{0.48\textwidth}
      \centering
      \includegraphics[width=\linewidth,height=1.8cm,keepaspectratio]{val_mae_benchmark.pdf}
      \vskip 1pt
      {\tiny \textbf{Hình 5: Đường cong hội tụ Validation MAE giữa các mô hình qua 5 seeds}}
    \end{column}
    \begin{column}{0.48\textwidth}
      \vskip 2pt
      \badge{ptitgreen}{\textbf{Độ chính xác số 1:}} Overall MAE đạt \textbf{3.1788}, tốt nhất toàn bảng.\\
      \badge{ptitblue}{\textbf{Ưu thế tầm nhìn xa:}} Tại $t+6$ (30 phút), đạt \textbf{3.4947} vs STAEformer (3.5120).\\
      \badge{ptitred}{\textbf{Hiệu quả phần cứng:}} Nhẹ hơn \textbf{70.4\% tham số} so với STAEformer.
    \end{column}
  \end{columns}
\end{frame}

% ------------------------------------------------------------------------------
% SLIDE 11: PHÂN TÍCH ĐỘNG HỌC SAI SỐ THEO HORIZON
% ------------------------------------------------------------------------------
\begin{frame}{10. Phân tích động học sai số theo Horizon ($t+1 \dots t+6$)}
  \begin{columns}[T]
    \begin{column}{0.53\textwidth}
      \centering
      \includegraphics[width=\linewidth,height=4.6cm,keepaspectratio]{error_by_horizon.pdf}
      \vskip 2pt
      {\tiny \textbf{Hình 6: Đường cong suy thoái sai số từ bước $t+1$ (5 phút) đến $t+6$ (30 phút)}}
    \end{column}

    \begin{column}{0.44\textwidth}
      \begin{block}{\textbf{Bảng MAE chi tiết qua từng bước nhìn xa}}
        \centering
        \scriptsize
        \renewcommand{\arraystretch}{1.08}
        \setlength{\tabcolsep}{3.5pt}
        \begin{tabular}{lcccc}
          \toprule
          \textbf{Horizon} & \textbf{STGCN} & \textbf{GWNet} & \textbf{STAEf.} & \textbf{TA-STGCN} \\
          \midrule
          $t+1$ (5 ph) & 2.9461 & 2.8715 & \textbf{2.8410} & 2.8524 \\
          $t+2$ (10 ph) & 3.1420 & 3.0540 & 3.0210 & \textbf{3.0180} \\
          $t+3$ (15 ph) & 3.3248 & 3.2234 & 3.1950 & \textbf{3.1892} \\
          $t+4$ (20 ph) & 3.4410 & 3.3610 & 3.3120 & \textbf{3.3015} \\
          $t+5$ (25 ph) & 3.5510 & 3.4720 & 3.4210 & \textbf{3.4080} \\
          $t+6$ (30 ph) & 3.6521 & 3.5800 & 3.5120 & \textbf{3.4947} \\
          \bottomrule
        \end{tabular}
      \end{block}
      \vskip 2pt
      \begin{beamercolorbox}[rounded=true,wd=\linewidth,sep=3pt]{block body example}
        \scriptsize
        \textbf{Nhận xét then chốt:} TA-STGCN có độ dốc tăng sai số phẳng nhất ($+0.642$ MAE từ $t+1 \to t+6$), kìm hãm đà phân kỳ sai số vượt trội so với các baseline.
      \end{beamercolorbox}
    \end{column}
  \end{columns}
\end{frame}

% ------------------------------------------------------------------------------
% SLIDE 12: ĐỘ BỀN VỮNG TRƯỚC NHIỄU CẢM BIẾN
% ------------------------------------------------------------------------------
\begin{frame}{11. Nghiên cứu độ bền vững trước nhiễu cảm biến camera}
  \begin{columns}[T]
    \begin{column}{0.54\textwidth}
      \centering
      \includegraphics[width=\linewidth,height=4.6cm,keepaspectratio]{noise_robustness_study.pdf}
      \vskip 2pt
      {\tiny \textbf{Hình 7: Mức tăng sai số MAE dưới tác động của nhiễu cảm biến (0\% -- 20\%)}}
    \end{column}

    \begin{column}{0.44\textwidth}
      \begin{alertblock}{\textbf{Khả năng kìm hãm sai số thực nghiệm}}
        \scriptsize
        \begin{itemize}
          \item \textbf{Mô hình không Attention (STGCN, GWNet):} Sai số tăng vọt (MAE lên $4.52$ và $4.32$ tại mức nhiễu $20\%$).
          \item \textbf{TA-STGCN:} Duy trì độ ổn định vượt bậc, kìm hãm MAE ở mức \textbf{3.7820} (thấp nhất trong tất cả các mô hình).
        \end{itemize}
      \end{alertblock}

      \begin{block}{\textbf{Kiểm định thống kê Wilcoxon \& Cohen's $d_z$}}
        \centering
        \scriptsize
        \renewcommand{\arraystretch}{1.05}
        \setlength{\tabcolsep}{4pt}
        \begin{tabular}{lccc}
          \toprule
          \textbf{Đối sánh cặp} & \textbf{$p$-value} & \textbf{Cohen's $d_z$} & \textbf{Thắng/Tổng} \\
          \midrule
          vs STGCN & $0.0039$ & \textbf{-2.84} & 5 / 5 seeds \\
          vs GWNet & $0.0156$ & \textbf{-1.98} & 5 / 5 seeds \\
          vs STAEformer & $0.0432$ & \textbf{-1.15} & 4 / 5 seeds \\
          \bottomrule
        \end{tabular}
      \end{block}
    \end{column}
  \end{columns}
\end{frame}

% ------------------------------------------------------------------------------
% SLIDE 13: MÔ HÌNH TRIỂN KHAI PHÂN TÁN EDGE - CLOUD
% ------------------------------------------------------------------------------
\begin{frame}{12. Thảo luận: Kiến trúc triển khai Edge--Cloud \& Tài nguyên tính toán}
  \begin{columns}[c]
    \begin{column}{0.49\textwidth}
      \centering
      \includegraphics[width=\linewidth,height=4.3cm,keepaspectratio]{edge_cloud_architecture.pdf}
      \vskip 2pt
      {\tiny \textbf{Hình 8: Phác thảo thiết kế ý niệm triển khai hệ thống phân tán Edge--Cloud}}
    \end{column}

    \begin{column}{0.49\textwidth}
      \begin{block}{\textbf{Phân bổ trách nhiệm tính toán phân tán}}
        \scriptsize
        \begin{itemize}
          \item \textbf{Tầng Biên (Edge Camera -- GĐ 1):} Chạy EfficientNet-B5 INT8 trích xuất vector $[\hat{c}_{\text{car}}, \hat{c}_{\text{bike}}]^T$. Hủy ảnh thô tại chỗ $\implies$ Bảo đảm \textbf{100\% quyền riêng tư} (Privacy-by-Design).
          \item \textbf{Tầng Trung tâm (Cloud Server -- GĐ 2):} Tiếp nhận vector số đếm nhỏ gọn, chạy TA-STGCN dự báo toàn mạng 608 nút.
        \end{itemize}
      \end{block}

      \begin{block}{\textbf{Hiệu quả tài nguyên phần cứng thực tế}}
        \scriptsize
        \begin{itemize}
          \item \textbf{Suy luận toàn mạng 608 nút:} Chỉ \textbf{1.23 ms} (nhanh gấp 15 lần so với STAEformer 18.60 ms).
          \item \textbf{VRAM GPU đào tạo:} 8.6 GB (so với STAEformer 11.5 GB), vừa vặn cấu hình phần cứng phổ thông.
        \end{itemize}
      \end{block}
    \end{column}
  \end{columns}
\end{frame}

% ------------------------------------------------------------------------------
% SLIDE 14: TÍNH MỚI, SÁNG TẠO & ĐÓNG GÓP THỰC TIỄN
% ------------------------------------------------------------------------------
\begin{frame}{13. Tính mới, Sáng tạo và Đóng góp thực tiễn}
  \begin{columns}[T]
    \begin{column}{0.48\textwidth}
      \begin{block}{\textbf{Tính mới và Đóng góp khoa học}}
        \scriptsize
        \begin{itemize}
          \item \textbf{Khung mô hình nối tầng phân tách:} Giải quyết toàn diện bài toán từ ảnh thô camera đường phố đến dự báo mạng lưới giao thông diện rộng.
          \item \textbf{Hồi quy đếm xe trực tiếp:} Xóa bỏ bẫy tính toán NMS, giải quyết hiệu quả bài toán che khuất tương hỗ của xe máy mật độ cao.
          \item \textbf{Cơ chế Temporal Attention khử nhiễu:} Chứng minh giải tích và thực nghiệm khả năng hoạt động như bộ lọc thông thấp thích ứng dữ liệu.
        \end{itemize}
      \end{block}
    \end{column}

    \begin{column}{0.48\textwidth}
      \begin{exampleblock}{\textbf{Ý nghĩa thực tiễn và Định hướng ứng dụng}}
        \scriptsize
        \begin{itemize}
          \item \textbf{Tận dụng hạ tầng camera sẵn có:} Không cần đầu tư xẻ đường chôn vòng từ cảm ứng tốn kém.
          \item \textbf{Bảo vệ quyền riêng tư công dân:} Hình ảnh người dân được xử lý và hủy tại biên, không truyền ảnh thô qua mạng công cộng.
          \item \textbf{Hỗ trợ điều hành giao thông thông minh (ITS):} Cung cấp dự báo sớm 15--30 phút giúp trung tâm điều phối tín hiệu đèn đón đầu tắc nghẽn.
        \end{itemize}
      \end{exampleblock}
    \end{column}
  \end{columns}
  \vskip 4pt
  \centering
  \badge{ptitgreen}{\textbf{Sản phẩm đề tài: Bộ mã nguồn hoàn chỉnh, Dataset IC4SD-Traffic-HCM \& Bản thảo báo cáo}}
\end{frame}

% ------------------------------------------------------------------------------
% SLIDE 15: KẾT LUẬN & ĐỊNH HƯỚNG PHÁT TRIỂN
% ------------------------------------------------------------------------------
\begin{frame}{14. Kết luận và Hướng mở rộng tiếp nối (Future Work)}
  \begin{columns}[T]
    \begin{column}{0.48\textwidth}
      \begin{block}{\textbf{Kết luận đề tài}}
        \scriptsize
        \begin{enumerate}
          \item \textbf{Hoàn thành toàn diện các mục tiêu:} Xây dựng thành công chuỗi mô hình hoàn chỉnh từ cảm nhận thị giác đến dự báo đồ thị không -- thời gian.
          \item \textbf{TA-STGCN đạt hiệu năng vượt bậc:} MAE tổng thể đạt 3.1788, kìm hãm xuất sắc sai số tại tầm nhìn xa 30 phút ($t+6$).
          \item \textbf{Hiệu quả tính toán:} Tiết kiệm 70.4\% tham số so với Transformer lớn nhất, suy luận thời gian thực 1.23 ms.
        \end{enumerate}
      \end{block}
    \end{column}

    \begin{column}{0.48\textwidth}
      \begin{exampleblock}{\textbf{Giới hạn và Hướng mở rộng tiếp nối}}
        \scriptsize
        \begin{itemize}
          \item \textbf{Mở rộng dữ liệu bán giám sát:} Khai thác Vision Foundation Models để tự động gán nhãn ảnh camera trong điều kiện ban đêm, thời tiết xấu.
          \item \textbf{Đồ thị không gian động:} Bổ sung đồ thị thích ứng tự học để nắm bắt các biến động tạm thời do phân luồng hoặc tai nạn.
          \item \textbf{Tích hợp yếu tố ngoại sinh:} Đưa dữ liệu mưa ngập đô thị, triều cường vào làm giàu đặc trưng đầu vào.
          \item \textbf{Triển khai thử nghiệm:} Kết nối điều khiển pha đèn tín hiệu thích ứng tại một số nút giao trọng điểm.
        \end{itemize}
      \end{exampleblock}
    \end{column}
  \end{columns}
  \vskip 6pt
  \centering
  \badge{ptitblue}{\textbf{Đề tài đã hoàn thành xuất sắc các nội dung nghiên cứu và cam kết khoa học}}
\end{frame}

% ------------------------------------------------------------------------------
% SLIDE 16: LỜI CẢM ƠN & Q&A
% ------------------------------------------------------------------------------
\begin{frame}[plain]
  \begin{tikzpicture}[remember picture,overlay]
    \fill[ptitblue] (current page.south west) rectangle (current page.north east);
    \node[anchor=north west] at ($(current page.north west)+(0.8cm,-0.4cm)$) {
      \includegraphics[height=1.0cm]{image2.png}
    };
    \node[anchor=north west, text=white, align=left] at ($(current page.north west)+(2.2cm,-0.45cm)$) {
      \small\textbf{HỌC VIỆN CÔNG NGHỆ BƯU CHÍNH VIỄN THÔNG}\\
      \scriptsize\textbf{KHOA CÔNG NGHỆ THÔNG TIN 1}
    };

    \node[text width=0.85\paperwidth, align=center, text=white] at ($(current page.center)+(0, 0.4cm)$) {
      {\Huge \textbf{XIN TRÂN TRỌNG CẢM ƠN!}}\\[8pt]
      {\large \textbf{QUÝ THẦY CÔ TRONG HỘI ĐỒNG VÀ CÁC BẠN}}\\[12pt]
      {\small Kính chúc Quý Thầy Cô dồi dào sức khỏe và công tác tốt!}\\[20pt]
      \badge{white}{\textcolor{ptitred}{\large \textbf{PHIÊN HỎI ĐÁP -- Q\&A}}}
    };

    \node[text width=0.85\paperwidth, align=center, text=white!75] at ($(current page.south)+(0, 0.9cm)$) {
      \scriptsize Sinh viên thực hiện: \textbf{Lê Việt Anh} -- Email: \texttt{vietanh@ptit.edu.vn} -- Học viện Công nghệ Bưu chính Viễn thông
    };
  \end{tikzpicture}
\end{frame}

\end{document}